\section{扩展议题：引导与校准}
\subsection{Self-Instruct 的 prompt 工程}
为了让 Self-Instruct 生成的 prompt 同时覆盖数学、代码、逻辑判断和对话，Weston 建议在 prompt 中明确写出 ``show your chain of thought''、``highlight where you doubt'' 和 ``produce a verification statement''。具体做法包括：
\begin{itemize}
    \item 使用模板：``Step-by-step, explain ..., then verify with a judgement block.''。
    \item 在 prompt 末尾加上 ``If you are unsure, state your uncertainty and ask for a cross-check''。
    \item 预设 verified response 的格式（reasoning + final verdict），让 reward model 训练时更容易匹配。
\end{itemize}
\mbox{}

\subsection{Calibration 与奖励归一}
Meta judge 训练需要 calibration，以便在 verified response 与 draft 之间公平判断。Weston 提醒我们 ``self evaluation is only going to help self-improve the model if inference compute for evaluation is controlled''（SRT 3236-3252）。常见策略包括 reward clipping、temperature scaling、cross-check rate scheduling 和 judge 再训练。

\begin{knowledgebox}{Calibration best practices}
对 judge 的 logits 做 temperature scaling，或在 verified response 出现低 confidence 时拉高 cross-check 频率；对 reward 值做 clipping 和 normalization，避免 single outlier 导致 self-rewarding 全面失控。
\end{knowledgebox}

\subsection{本章小结}
Prompt 设计与 calibration 是 self-rewarding 成功的必要条件，前者确保任务覆盖而不失方向感，后者则让 judge 在 cross-check 中更稳定、不会因为 hallucination 自信而偏离。

